\documentclass[11pt]{article}

\usepackage[margin=1in]{geometry}
\usepackage{amsmath, amssymb}
\usepackage{xcolor}
\usepackage{listings}
\usepackage{hyperref}
\usepackage{array}
\usepackage{booktabs}

\hypersetup{
    colorlinks=true,
    linkcolor=blue!60!black,
    urlcolor=blue!60!black,
    citecolor=blue!60!black
}

\lstdefinestyle{bugpython}{
  language=Python,
  basicstyle=\ttfamily\small,
  keywordstyle=\color{blue!60!black},
  stringstyle=\color{green!40!black},
  commentstyle=\color{gray!70!black},
  showstringspaces=false,
  breaklines=true,
  frame=single,
  rulecolor=\color{gray!30},
  columns=fullflexible,
  keepspaces=true
}

\title{SymPy Bug Report}
\author{}
\date{}

\begin{document}

\maketitle

\section{Bug 1: solveset returns 0 for asin outside the principal range}

\subsection{Minimal Reproducer}

\medskip

\begin{lstlisting}[style=bugpython]
import sympy
from sympy import Eq, S, asin, pi, simplify, solveset, symbols

x = symbols("x")
eq = Eq(asin(x), pi)
sol = solveset(eq, x, domain=S.Complexes)

print("SymPy version:", sympy.__version__)
print("solveset result:", sol)
print("residual at returned value 0:",
      simplify(eq.lhs.subs(x, 0) - eq.rhs.subs(x, 0)))
\end{lstlisting}

\subsection{Observed Behavior}

\medskip

\begin{lstlisting}[style=bugpython]
SymPy version: 1.14.0
solveset result: {0}
residual at returned value 0: -pi
\end{lstlisting}

SymPy returns the finite solution set \(\{0\}\).  However, direct substitution
of the returned value into the original equation gives
\[
  \operatorname{asin}(0)-\pi = -\pi \ne 0.
\]
Thus the returned element does not satisfy the equation that \texttt{solveset}
was asked to solve.

\subsection{Expected Behavior}

The mathematically correct result is \(\varnothing\), represented in SymPy as
\texttt{EmptySet}.  There is no complex value \(x\) for which SymPy's principal
inverse sine satisfies \(\operatorname{asin}(x)=\pi\).

\subsection{Mathematical Explanation}

SymPy's \(\operatorname{asin}\) is the principal inverse sine.  On real inputs its image is
\([-\pi/2,\pi/2]\), and more generally the complex principal branch has a fixed
principal image rather than all complex values.  The target \(\pi\) is outside
that principal image.

Applying \(\sin\) to both sides gives \(x=\sin(\pi)=0\), but this implication
is not reversible for the principal inverse sine.  The value \(x=0\) satisfies
\(\sin(\pi)=0\), but it does not satisfy the original equation because
\(\operatorname{asin}(0)=0\), not \(\pi\).  The wrong answer is therefore not a harmless
branch choice, a missing constant, or an equivalent representation of a
solution; it is a concrete non-solution.

\subsection{Source-Level Diagnosis}

The diagnosis identifies the relevant code path in
\nolinkurl{sympy/solvers/solveset.py}, in the complex inversion helper
\texttt{\_invert\_complex} near line 550.  The traced call path is that
\texttt{solveset} converts \texttt{Eq(asin(x), pi)} to
\texttt{asin(x) - pi = 0} and enters \texttt{\_solveset}.  The expression is
not handled by \texttt{\_solve\_trig}, because the top-level function is the
inverse trigonometric function \texttt{asin}, not \texttt{sin}.  The solver then
calls \texttt{\_invert(..., domain=S.Complexes)}, which dispatches to
\texttt{\_invert\_complex}.

For this equation, \texttt{\_invert\_complex(asin(x) - pi, \{0\}, x)} removes
the additive constant and then handles \(\operatorname{asin}(x)\) through the generic
function-inverse path.  It applies \texttt{asin.inverse()}, i.e. \(\sin\), to
the right-hand side and produces \((x,\{\sin(\pi)\})=(x,\{0\})\).  The
problematic source rule is the generic use of a function's inverse when one is
available:

\begin{lstlisting}[style=bugpython]
return _invert_complex(f.args[0],
                       imageset(Lambda(n, f.inverse()(n)), g_ys), symbol)
\end{lstlisting}

For principal \(\operatorname{asin}\), this treats \(\sin\) as a globally valid inverse
without carrying the required range condition on the right-hand side.  The
later filtering accepts the finite inversion result after \texttt{domain\_check}
rather than checking the residual \(\operatorname{asin}(0)-\pi\) against the original
equation.  Consequently, the candidate \(0\) is accepted even though it is not
a solution.  A plausible fix direction supported by the diagnosis is to make
inverse-function handling for principal inverse functions carry image or range
conditions, returning \(\varnothing\) or a \texttt{ConditionSet} when a concrete
target such as \(\pi\) violates the principal \(\operatorname{asin}\) image.

\subsection{Additional Failing Instances}

The same failure occurs for integer multiples of \(\pi\) that lie outside the
principal image of \(\operatorname{asin}\).  The independent verification checks
\[
  \operatorname{asin}(x)=n\pi,\qquad n\in\{1,2,-1,3,4,5\},
\]
over the complex numbers.  In each case, SymPy's inversion computes
\(x=\sin(n\pi)=0\), but the returned value has residual
\(\operatorname{asin}(0)-n\pi=-n\pi\ne0\).  Since none of these targets belongs to the
principal image of \(\operatorname{asin}\), the expected solution set is \(\varnothing\) in
all cases.

\begin{center}
\small
\renewcommand{\arraystretch}{1.3}
\begin{tabular}{@{}>{\raggedright\arraybackslash}p{0.44\linewidth}>{\raggedright\arraybackslash}p{0.23\linewidth}>{\raggedright\arraybackslash}p{0.23\linewidth}@{}}
\toprule
\textbf{Input / Expression} & \textbf{SymPy behavior} & \textbf{Expected behavior} \\
\midrule
\(\operatorname{solveset}(\operatorname{asin}(x)=\pi,\ x,\ \mathbb{C})\) & returns \(\{0\}\); residual \(-\pi\) & \(\varnothing\) \\
\(\operatorname{solveset}(\operatorname{asin}(x)=2\pi,\ x,\ \mathbb{C})\) & returns \(\{0\}\); residual \(-2\pi\) & \(\varnothing\) \\
\(\operatorname{solveset}(\operatorname{asin}(x)=-\pi,\ x,\ \mathbb{C})\) & returns \(\{0\}\); residual \(\pi\) & \(\varnothing\) \\
\(\operatorname{solveset}(\operatorname{asin}(x)=3\pi,\ x,\ \mathbb{C})\) & returns \(\{0\}\); residual \(-3\pi\) & \(\varnothing\) \\
\(\operatorname{solveset}(\operatorname{asin}(x)=4\pi,\ x,\ \mathbb{C})\) & returns \(\{0\}\); residual \(-4\pi\) & \(\varnothing\) \\
\(\operatorname{solveset}(\operatorname{asin}(x)=5\pi,\ x,\ \mathbb{C})\) & returns \(\{0\}\); residual \(-5\pi\) & \(\varnothing\) \\
\bottomrule
\end{tabular}
\end{center}

\subsection{Related Issues Assessment}

The bug-finding harness performed a best-effort deduplication check.  The
closest related public material found was documentation confirming the
principal-branch behavior of \(\operatorname{asin}\), general principal-branch references,
and SymPy documentation stating that \texttt{solveset} returns values for which
the equation is true.  No public issue, pull request, discussion, or external
report was identified for \(\operatorname{asin}(x)=\pi\), the wrong output \(\{0\}\), or a
missing principal-range check in \texttt{solveset} for this exact case.  The
deduplication verdict is therefore a new instance of a known failure mode: the
general need for branch- and range-aware inverse solving is known, but this
specific reproducible case appears previously unreported and remains
individually actionable as an issue and regression test.

\subsection{Proposed SymPy Regression Test}

\medskip

\begin{lstlisting}[style=bugpython]
import pytest

from sympy import Eq, FiniteSet, S, asin, pi, simplify, solveset, symbols


@pytest.mark.parametrize("n", [1, 2, -1, 3, 4, 5])
def test_solveset_principal_asin_rejects_out_of_range_targets(n):
    x = symbols("x")
    eq = Eq(asin(x), n*pi)
    sol = solveset(eq, x, domain=S.Complexes)

    assert sol != FiniteSet(0)
    assert simplify(eq.lhs.subs(x, 0) - eq.rhs.subs(x, 0)) != 0
\end{lstlisting}

\end{document}
